\subsection{\problemblock{*File} Data Block}\label{sec:trajectory-block-file}

The trajectory-level \problemblock{*file} block creates a file containing columns
of trajectory information. Each selected output variable is written as a column of
numbers separated by at least one space. Column headings appear once at the
beginning, and the columns are continuous thereafter. This format is well suited for
import into plotting programs, spreadsheets, or other analysis software because it
contains no page headings or page breaks; at most, the initial column-heading line
must be removed.

\begin{lstlisting}[style=taosmanual]
*file traj.dat
  time alt long latgd range vel
\end{lstlisting}

writes the selected trajectory variables to \texttt{traj.dat}. The beginning of such
a file is shown in \cref{fig:trajectory-output-file-example}.

\begin{figure}[htbp]
  \centering
  \begin{minipage}{0.93\linewidth}
    \lstinputlisting[style=taosmanual,basicstyle=\ttfamily\scriptsize]{examples/chapter04/trajectory-output-example.txt}
  \end{minipage}
  \caption{Trajectory output file example.}
  \label{fig:trajectory-output-file-example}
\end{figure}

The filename must be valid for the operating system. If a file already exists, TAOS
overwrites it without warning and the previous data are lost. This is normally
acceptable because the file is intended to contain the most recent trajectory run.

Any standard output variable or user-defined output variable may be written. A
user-defined variable referenced in a \problemblock{*file} block must have been
defined earlier. Units and output format are controlled by
\problemblock{*units/fmt}.

Each output line is limited to 400 characters, which limits the list to approximately
30 variables. In practice, 12 or fewer is recommended because some screen editors
and printers handle only about 130 characters per line.

Radar and relative-vehicle variables, whose names begin with \texttt{rad} or
\texttt{rel}, require a square-bracket subscript. For example,
\statevar{radrng[2]} is the current trajectory's range from radar station 2, and
\statevar{relvel[4]} is the closing velocity of trajectory 4 relative to the current
trajectory.

Any number of trajectory-level \problemblock{*file} blocks can be included. A
separate output file is created for each block.
